\needspace{12\baselineskip}
\SourceLeaf{319}
\pdfbookmark[1]{16. On the Zodiacs of Celestial Bodies}{zusatz-016}
\begin{center}
[On the Zodiacs of Celestial Bodies.]
\end{center}

\begin{center}
[16.]
\end{center}

\par\noindent\rule{\textwidth}{.3pt}\par

\noindent\small Astronomische Nachrichten. Vol. XXVII. No. 625. Page 1. January 29, 1848.\normalsize

\par\noindent\rule{\textwidth}{.3pt}\par

\begin{center}
Letter from Privy Court Councillor \textsc{Gauss} to the Editor.

Göttingen, November 23, 1847.
\end{center}

The second set of elements of Iris, published in No. 615 of the A.\,N. and calculated by Prof. \textsc{Goldschmidt} immediately after September 19\textsuperscript{th}, continued to agree well with the observations in the following two months (on November 17\textsuperscript{th} and 18\textsuperscript{th} the difference from Mr. \textsc{Rümker}'s meridian observation was \(16^{\prime\prime}\) in R.A.\ and \(2^{\prime\prime}\) in declination). We are therefore justified in regarding these elements as already a very close approximation, and I accordingly asked Prof. \textsc{Goldschmidt} to use them to calculate the zodiac of Iris. I give the result of this work below, as it may help in searching earlier observations for the same object.

In taking up again on this occasion the paper\HistoricalNoteEight{Vol. VI, p. 106.]} in which, 44 years ago, I gave the general method of determining the boundaries of such a zodiac, I see that I already indicated the exceptional cases in which the region of the celestial sphere where a celestial body orbiting the Sun according to \textsc{Kepler}'s laws can appear geocentrically has either no boundary at all or only one, although in the former case the method gives one closed
\SourceLeaf{320}curve and in the latter two such curves. I also raised the question there of what the curves obtained by calculation actually mean in such exceptional cases. At the time I confined myself to these hints, since a fuller treatment would have been a hors d’oeuvre there, and I also wished to leave others the pleasure of working on a mathematical problem that I considered far from uninteresting. In the long interval since then the question has entirely slipped my mind, but as I have heard of no one else taking it up, I shall use this opportunity to give at least the essentials needed to answer it.

A geocentric position of the planet (or comet) in question is obtained by combining a point of its orbit with a point of the Earth's orbit, but one and the same geocentric position can also arise from two, three, or four different combinations. To make the construction clear, pass a plane through the Sun perpendicular to the straight line representing a given geocentric position on the celestial sphere, and project both the Earth's orbit and the planet's orbit orthographically onto that plane. Both projections are ellipses, or more generally conic sections, which will intersect either nowhere, twice, or four times. A tangency is to be regarded as the coalescence of two intersections. Each intersection represents a geocentric position that either coincides with the given position or lies opposite to it on the celestial sphere. This depends on whether the point of the planet's orbit or the point of the Earth's orbit is higher, the two points having the same projection. The upper side of the plane of projection is taken to be the side on which the given geocentric position lies.

It follows that an arbitrary point on the celestial sphere, taken as a possible geocentric position, can correspond to no combination at all, to one combination, or to two, three, or four combinations. For a given planet, the entire surface of the celestial sphere is thus divided into regions. The curves determined by my method are, generally speaking, simply the dividing curves between two regions of the sphere, where the points on one side can be geocentric positions in two more ways than those on the
\SourceLeaf{321}other side. The points on the dividing curves mark the transition, that is, they can be geocentric positions in one more way than the points on one side and in one fewer way than the points on the other. One can also give a criterion for deciding a priori on which side of the dividing curve there are two more solutions than on the other, a matter on which I will not dwell at present\HistoricalNoteEight{See Note 17.]}.

All the planets known thus far have orbits such that the dividing curves found by the theory are always actual boundaries. There are two dividing curves that partition the celestial sphere into three regions. Two are isolated regions into which no geocentric positions fall at all. The third, a belt between them, contains all the geocentric positions. Each point within the belt can be a geocentric position in two ways and each point on the boundary in one way.

One can, however, imagine hypothetical orbits\HistoricalNoteDouble{See Note 18.]} (several cometary orbits may fall into this case, though I have made no investigation of the matter so far) for which two boundaries again divide the celestial sphere into three parts. One part contains no geocentric positions at all, another contains the geocentric positions that occur in two ways, and the third contains the points that can be geocentric positions in four different ways.

There is still more variety if one dividing curve crosses itself once or twice, forming a single or a double loop. In the first case it will cut off two regions from the belt-shaped part, and in the second case three, in which there are 4 and 0 or 4, 0, 4 solutions, respectively.

An orbit that interlocks with the Earth's orbit like a link in a chain is different again. In such a case there is only one connected boundary curve, and thus two distinct regions. The points in one part can be geocentric positions in one way and those in the other in three ways.

\needspace{29\baselineskip}
\SourceLeaf{322}
\begin{center}
Zodiac of Iris.
\end{center}
\begingroup\small
\setlength{\tabcolsep}{7pt}
\renewcommand{\arraystretch}{1.15}
\begin{center}
\begin{tabular}{|r|r@{\;}r|r@{\;}r|}
\multicolumn{1}{|c|}{Right}&\multicolumn{4}{c|}{Declination}\\
\multicolumn{1}{|c|}{Asc.}&\multicolumn{2}{c|}{northern boundary}&\multicolumn{2}{c|}{southern boundary}\\
\(0^{0}\)&\(+12^{0}\)&\(9^{\prime}\)&\(+3^{0}\)&\(44^{\prime}\)\\
5&14&10&5&41\\
10&16&5&7&33\\
15&17&51&9&19\\
20&19&30&10&57\\
25&21&0&12&29\\
30&22&21&13&52\\
35&23&33&15&7\\
40&24&36&16&14\\
45&25&29&17&11\\
50&26&12&17&59\\
55&26&45&18&39\\
60&27&9&19&8\\
65&27&24&19&29\\
70&27&29&19&41\\
75&27&24&19&42\\
80&27&11&19&36\\
85&26&48&19&19\\
90&26&15&18&54\\
95&25&33&18&19\\
100&24&41&17&35\\
105&23&41&16&42\\
110&22&32&15&41\\
115&21&14&14&31\\
120&19&47&13&12\\\hline
\end{tabular}
\end{center}
\begin{center}
\begin{tabular}{|r|r@{\;}r|r@{\;}r|}
\multicolumn{1}{|c|}{Right}&\multicolumn{4}{c|}{Declination}\\
\multicolumn{1}{|c|}{Asc.}&\multicolumn{2}{c|}{northern boundary}&\multicolumn{2}{c|}{southern boundary}\\
\(125^{0}\)&\(+18^{0}\)&\(12^{\prime}\)&\(+11^{0}\)&\(46^{\prime}\)\\
130&16&29&10&12\\
135&14&40&8&31\\
140&12&45&6&45\\
145&10&44&4&53\\
150&8&39&2&56\\
155&6&31&\(+0\)&56\\
160&4&22&\(-1\)&6\\
165&\(+2\)&12&3&8\\
170&0&0&5&11\\
175&\(-2\)&9&7&13\\
180&4&14&9&11\\
185&6&16&11&6\\
190&8&12&12&56\\
195&10&2&14&40\\
200&11&45&16&19\\
205&13&20&17&50\\
210&14&46&19&13\\
215&16&4&20&58\\
220&17&13&21&35\\
225&18&12&22&33\\
230&19&1&23&23\\
235&19&40&24&4\\
240&20&10&24&35\\\hline
\end{tabular}
\end{center}
\begin{center}
\begin{tabular}{|r|r@{\;}r|r@{\;}r|}
\multicolumn{1}{|c|}{Right}&\multicolumn{4}{c|}{Declination}\\
\multicolumn{1}{|c|}{Asc.}&\multicolumn{2}{c|}{northern boundary}&\multicolumn{2}{c|}{southern boundary}\\
\(245^{0}\)&\(-20^{0}\)&\(29^{\prime}\)&\(-24^{0}\)&\(58^{\prime}\)\\
250&20&38&25&11\\
255&20&38&25&15\\
260&20&27&25&11\\
265&20&5&24&56\\
270&19&34&24&33\\
275&18&52&24&1\\
280&18&0&23&19\\
285&16&58&22&28\\
290&15&46&21&28\\
295&14&25&20&20\\
300&12&54&19&2\\
305&11&14&17&36\\
310&9&26&16&2\\
315&7&30&14&21\\
320&5&28&12&34\\
325&3&21&10&40\\
330&\(-1\)&9&8&41\\
335&\(+1\)&6&6&39\\
340&3&22&4&34\\
345&5&37&2&28\\
350&7&51&\(-0\)&22\\
355&10&2&\(+1\)&42\\
360&12&9&3&44\\\hline
\end{tabular}
\end{center}
\endgroup
\par\medskip\begin{center}\rule{.22\textwidth}{.3pt}\end{center}
